\chapter{ДҮГНЭЛТ}

\section{Судалгааны үр дүнгийн нэгтгэл}

Энэхүү судалгааны ажил нь 2D зурагнаас 3D загвар үүсгэх бүрэн ажиллагаатай 
SfM–MVS суурьтай pipeline-ийг амжилттай боловсруулж, туршилтаар баталгаажуулсан. 
Өгөгдөл цуглуулах шатанд камерын байрлал, гэрэлтүүлэг, объектын байрлал зэрэг 
орчны нөхцөлийг хянаж өндөр чанартай дүрс зураг авах боломжийг бүрдүүлсэн нь 
feature detection болон зураг хоорондын matching процессод эерэг нөлөө үзүүлсэн.

\subsection{Техникийн үр дүн}

4-р бүлэгт дэлгэрэнгүй авч үзсэн туршилтын гол үзүүлэлтүүдийг дараах хүснэгтэд нэгтгэн харуулав.

\begin{table}[H]
\centering
\caption{Системийн гол гүйцэтгэлийн үзүүлэлтүүд}
\begin{tabular}{|l|r|}
\hline
\textbf{Үзүүлэлт} & \textbf{Утга} \\
\hline
Туршилтын датасет (зураг) & 47 \\
Sparse point cloud & 35{,}000+ цэг \\
Dense point cloud & 2.1 сая цэг \\
Mesh vertices / faces & 1.2 сая / 2.3 сая \\
Хамгийн сайн Chamfer Distance (1B: LightGlue+Exhaustive) & 0.0074 м \\
Хамгийн сайн RMSE (1A: SIFT+Exhaustive) & 0.0064 м \\
\hline
\end{tabular}
\label{tab:key_results}
\end{table}

SfM шатанд sparse point cloud амжилттай үүсгэж, зураг хоорондын pose estimation-ийг баталгаажуулсан. MVS шатанд depth map фьюзионоор нягт цэгэн үүл гаргаж, Poisson reconstruction ашиглан mesh үүсгэж текстур зураглал хийсэн.

\subsection{Төхөөрөмжийн үр дүн}

Туршилтын pipeline-г i9 CPU, RTX 4060 GPU, 16GB RAM бүхий машин дээр ажиллуулж, Python, PyQt5, COLMAP, OpenMVS сан ашиглав. Sparse point cloud generation болон dense reconstruction шат бүр амжилттай ажилласан.

\textbf{Гол үзүүлэлтүүд:}
\begin{itemize}
    \item Sparse point cloud: 35,000+ цэг
    \item Dense point cloud: 2.1 сая цэг
    \item Mesh vertices: 1.2 сая, faces: 2.3 сая
    \item Хамгийн сайн Chamfer Distance: \textbf{0.0074~м} (1B: LightGlue + Exhaustive)
    \item Хамгийн сайн RMSE: \textbf{0.0064~м} (1A: SIFT + Exhaustive)
\end{itemize}

\subsection{Зорилтын биелэлт}

1-р бүлэгт тавьсан зорилтуудын биелэлтийг дараах хүснэгтэд харуулав.

\begin{table}[H]
\centering
\caption{Зорилтын биелэлт}
\begin{tabular}{|c|p{9cm}|c|}
\hline
\textbf{№} & \textbf{Зорилт} & \textbf{Биелэлт} \\
\hline
1 & Controlled lighting box зохион бүтээж, тогтвортой нөхцөлд зураг авах орчин бий болгох & \checkmark \\
\hline
2 & Объектуудыг олон өнцгөөс жигд гэрэлтүүлэгтэйгээр зураг авах & \checkmark \\
\hline
3 & COLMAP ашиглан SfM ба MVS реконструкц хэрэгжүүлэх & \checkmark \\
\hline
4 & PyQt5-д суурилсан GUI систем бүтээх & \checkmark \\
\hline
5 & Chamfer Distance болон RMSE-ээр гүйцэтгэлийг үнэлэх & \checkmark \\
\hline
6 & SfM → MVS → Mesh алхмуудыг автоматаар ажиллуулдаг 9 үе шаттай pipeline боловсруулах & \checkmark \\
\hline
\end{tabular}
\label{tab:objectives_completion}
\end{table}

% \section{Зорилтын биелэлт}

% \begin{table}[H]
% \centering
% \caption{Зорилтын биелэлт}
% \begin{tabular}{|c|l|c|l|}
% \hline
% \textbf{№} & \textbf{Зорилт} & \textbf{Биелэлт} & \textbf{Тайлбар} \\
% \hline
% 1 & Өгөгдөл цуглуулах & \checkmark & 40-80 зураг, controlled lighting \\
% \hline
% 2 & Pre-processing хийх & \checkmark & Resize, denoise, color normalization \\
% \hline
% 3 & Feature detection & \checkmark & SIFT + SuperPoint, BFMatcher/FLANN \\
% \hline
% 4 & Camera calibration & \checkmark & Checkerboard ашигласан \\
% \hline
% 5 & SfM & \checkmark & Sparse point cloud үүсгэсэн \\
% \hline
% 6 & MVS & \checkmark & Dense point cloud гаргасан \\
% \hline
% 7 & Mesh reconstruction & \checkmark & Poisson reconstruction, texture baking \\
% \hline
% \end{tabular}
% \label{tab:objectives}
% \end{table}

\section{Шинэлэг хувь нэмэр}

\begin{enumerate}
    \item \textbf{PyQt5-д суурилсан GUI боловсруулсан} — бүх 9 үе шатыг нэг цонхноос удирдах боломжтой, лог харуулах, abort хийх боломж бүхий хэрэглэхэд хялбар интерфейс; энэ нь pipeline-ийг мэргэжлийн бус хэрэглэгчдэд ч хүртээмжтэй болгосон
    \item \textbf{Matching стратегийн нарийвчлалд үзүүлэх нөлөөг туршилтаар тодорхойлсон} — гурван ялгаатай туршилтын цуваа явуулж хослолын арга (exhaustive vs sequential), overlap параметр болон зургийн тооны нарийвчлалд үзүүлэх нөлөөг тодорхойлсон; sequential matching дэхь overlap=10 нь оновчтой босго утга болох нь тогтоогдсон
    \item \textbf{Chamfer Distance болон RMSE-д суурилсан тоон үнэлгээний аргачлалыг хэрэглэсэн} — зөвхөн харааны үнэлгээнд тулгуурлахгүйгээр объектив метрик ашиглан реконструкцийн чанарыг хэмжсэн
\end{enumerate}

\section{Хязгаарлалт}

\begin{enumerate}
    \item Зөвхөн нэг объектын туршилт хийсэн бөгөөд олон объектод ажиллах боломж нэмэлт туршилт шаардана
    \item Тогтмол эргэдэг тавцан болон гэрлийн хайрцаг ашигласан тул гадна буюу динамик орчинд ажиллах тохиолдлыг хамраагүй
    \item GUI-ийн хэрэглэгчийн туршлага (UX) цаашид сайжруулах шаардлагатай — камерын калибрацийн интерфейс, зургийн урьдчилсан боловсруулалтын сонголтууд дутмаг
    \item Sequential matching арга нь хязгаарлагдмал зургийн тооны датасетэд зарим matcher-д бүтэлгүйтэж болно. Туршилтаар 49 зурагтай датасет дэхь sequential matching нь LightGlue matcher-тэй хослуулахад Chamfer Distance 0.143~м (exhaustive тохиргооноос 19 дахин муу) гаргасан. Нейрон сүлжээнд суурилсан фичер илрүүлэлтийн арга ашиглах нь энэ асуудлыг шийдвэрлэх боломжтой боловч matching стратегийн сонголт ч адилхан чухал ач холбогдолтой болох нь тогтоогдлоо.
\end{enumerate}

\section{Цаашдын судалгааны чиглэл}

Боловсруулсан pipeline-ийн туршилтын явцад илэрсэн хязгаарлалтуудад тулгуурлан 
дараах чиглэлүүдэд судалгааг үргэлжлүүлэх нь зүйтэй гэж үзэж байна.

\begin{enumerate}
    \item \textbf{ML-д суурилсан фичер тохиролт нэвтрүүлэх} — Энэ судалгааны хамгийн 
    тулгамдсан асуудал нь COLMAP-ийн exhaustive matcher зарим нөхцөлд зургуудын 
    хоорондын тохиролтыг найдвартай хийж чадахгүй байсан явдал юм. SuperPoint болон
    SuperGlue нейрон сүлжээнд суурилсан фичер илрүүлэлт болон тохиролтын аргуудыг
    pipeline-д нэгтгэх нь энэ асуудлыг шийдвэрлэх хамгийн тохиромжтой арга зам юм. 
    SuperPoint нь масштаб болон гэрэлтүүлгийн өөрчлөлтөд SIFT-ээс илүү тэсвэртэй 
    байдаг бөгөөд SuperGlue нь зургуудын хооронд харагдах давхцал бага байх үед ч 
    найдвартай тохиролт хийх чадвартай. Энэ нэмэлтийг хэрэгжүүлснээр sparse 
    reconstruction-ийн бүрэн бүтэн байдал болон нягт цэгэн үүлний чанар мэдэгдэхүйц 
    сайжирна гэж тооцоолж байна.

    \item \textbf{Гүйцэтгэлийн оновчлол} — CUDA-р хурдасгах болон олон тэрдэлтэй 
    (multi-threaded) файл боловсруулалт нэмж нийт pipeline-ийн ажиллагааны цагийг 
    богиносгох. Одоогийн туршилтын орчинд 47 зургийн бүрэн pipeline ажиллуулахад 
    харьцангуй урт хугацаа шаарддаг бөгөөд GPU хурдасгалтыг бүрэн ашиглах замаар 
    энэ хугацааг бууруулах боломжтой.

    \item \textbf{GUI өргөтгөл} — Камерын калибрацийн интерфейс, зургийн урьдчилсан 
    боловсруулалтын тохиргооны сонголтууд болон дотоод 3D дүрслэгч нэмэх нь 
    хэрэглэгчийн туршлагыг улам сайжруулна.

    \item \textbf{NeRF болон Gaussian Splatting-тэй харьцуулалт} — Орчин үеийн гүн 
    сургалтад суурилсан аргуудтай тоон болон харааны харьцуулалт хийх нь 
    COLMAP+OpenMVS pipeline-ийн давуу болон сул талуудыг тодорхой болгоход тусална.
\end{enumerate}

\section{Төгсгөлийн үг}

Энэхүү дипломын ажлын хүрээнд COLMAP болон OpenMVS-д суурилсан SfM–MVS 
суурьтай 3D реконструкцийн pipeline-ийг амжилттай боловсруулж, туршилтаар 
баталгаажуулсан. Feature detection, sparse болон dense point cloud generation, 
mesh reconstruction, texture mapping гэсэн шат бүрийг автоматжуулсан 9 үе шаттай 
pipeline нь нэг цогц системд нэгтгэгдсэн. Үүнээс гадна, бүх үе шатыг нэг
цонхноос удирдах боломжтой PyQt5-д суурилсан GUI програмыг амжилттай
хэрэгжүүлсэн нь системийг мэргэжлийн бус хэрэглэгчдэд ч ашиглах боломжтой
болгосон гол ахиц юм.

Гурван ялгаатай туршилтын цуваа явуулж matching параметрүүдийн нарийвчлалд үзүүлэх нөлөөг тодорхойлсон бөгөөд хамгийн сайн Chamfer Distance нь 1B (LightGlue + Exhaustive) тохиргоонд \textbf{0.0074~м}, хамгийн сайн RMSE нь 1A (SIFT + Exhaustive) тохиргоонд \textbf{0.0064~м} гарсан нь системийн реконструкцийн нарийвчлал хангалттай өндөр байгааг харуулж байна. Цаашдын 
судалгаанд exhaustive matching-ийн хязгаарлалтыг SuperPoint/SuperGlue-д эсвэл DINOv2-д суурилсан 
ML фичер тохиролтоор орлуулах нь pipeline-ийн ерөнхий чанарыг мэдэгдэхүйц 
дээшлүүлэх боломжтой гол алхам болно.


